\subsection{到局部紧空间中的逆紧映射}

命题 7. 设 f 是豪斯多夫空间 X 到局部紧空间 Y 中的一个连续映射。则 f 逆紧当且仅当 Y 的每个紧子集在 f 下的逆像是紧的。此外，若 f 逆紧，则 X 是局部紧的。

若 \(f\) 逆紧且 \(\mathbf{K}\) 是 \(\mathbf{Y}\) 的一个紧子集，则由第 2 目命题 6，\(\overline{f}^{1}(\mathbf{K})\) 是紧的。反之，若这个条件满足，设 \((\mathbf{U}_{\alpha})\) 是 \(\mathbf{Y}\) 的一个由相对紧开集组成的覆盖。则诸集合 \(\overline{f}^{1}(\overline{\mathbf{U}}_{\alpha})\) 在 \(\mathbf{X}\) 中是紧的，且它们的内部覆盖 \(\mathbf{X}\)；由于 \(\mathbf{X}\) 是豪斯多夫的，这表明 \(\mathbf{X}\) 是局部紧的。此外，诸映射 \(f_{\overline{\mathbf{U}}_{\alpha}}: \overline{f}^{1}(\overline{\mathbf{U}}_{\alpha}) \to \overline{\mathbf{U}}_{\alpha}\) 中的每一个都是逆紧的（第 2 目，定理 1，推论 2），因此由第 1 目命题 3 b)，\(f\) 是逆紧的。

推论. 设 X, X' 是两个局部紧空间，Y（相应地，Y'）是由 X（相应地，X'）添加无穷远点 ω（相应地，ω′）而得到的紧空间（§ 9，第 8 目）。则连续映射 f: X → X' 逆紧，当且仅当它的延拓 \(\overline{f}\) : Y → Y′（满足 \(\overline{f}(\omega) = \omega'\)）是连续的。

由命题 7，\(f\) 逆紧当且仅当对每个紧子集 \(\mathbf{K}'\)，\(\mathbf{X}', \bar{f}^1(\mathbf{X}' - \mathbf{K}') = \mathbf{X} - \bar{f}^1(\mathbf{K}')\) 是 \(\mathbf{X}\) 的一个紧子集的余集；由 \(\omega\)（相应地，\(\omega'\)）在 \(\mathbf{Y}\)（相应地，\(\mathbf{Y}'\)）中邻域的定义，此事成立当且仅当 \(\bar{f}\) 在 \(\omega\) 处连续。

